% ==============================================================================
%  2.1  El análisis de regresión
% ==============================================================================
\section{El análisis de regresión}
\label{sec:t2-analisis}

El término \textbf{regresión} se utiliza para describir relaciones estadísticas entre
variables. El \textbf{análisis de regresión} es la técnica estadística que permite modelizar e
investigar la relación de una variable $Y$ con otra u otras variables $X_1,\ldots,X_k$
($Y \leftarrow X_1,\ldots,X_k$), es decir, estudiar cómo afectan a $Y$ los cambios en las $X_i$.
\begin{itemize}
  \item $Y$: \textbf{variable respuesta} o \textbf{dependiente}. Es una variable aleatoria y la
    de mayor interés en el problema.
  \item $X_1,\ldots,X_k$: \textbf{regresores}, \textbf{predictores}, variables
    \textbf{independientes} o \textbf{explicativas}. Aportan información sobre la variabilidad
    de $Y$, y el interés está en evaluar el efecto que sus cambios producen en $Y$.
\end{itemize}

\subsection{Componentes del modelo de regresión}

El modelo de regresión debe expresar formalmente las dos componentes de una relación
estadística (\cref{sec:t1-relacion}):
\begin{enumerate}
  \item la tendencia \textbf{sistemática} con la que cambia $Y$ al cambiar los predictores,
    $\vect{X}=(X_1,\ldots,X_k)$;
  \item la componente \textbf{aleatoria}, que recoge la variabilidad con la que los datos se
    distribuyen alrededor de esa tendencia.
\end{enumerate}
Para ello se supone que:
\begin{enumerate}
  \item existe una \textbf{distribución de $Y$ para cada nivel de $\vect{X}$};
  \item la \textbf{media} de estas distribuciones condicionadas \textbf{varía de forma
    sistemática con $\vect{X}$}.
\end{enumerate}
Los modelos son, por tanto, de la forma
\[
  Y = m(X_1,\ldots,X_k) + \varepsilon,
  \qquad
  m(x_1,\ldots,x_k) = \E\left(Y \cond X_1=x_1,\ldots,X_k=x_k\right),
\]
donde $m$ es la función de regresión (\cref{def:funcion-regresion}).

\begin{ejemplo}[Evaluación de empleados][ej:t2-empleados]
  Se quiere evaluar la relación entre la eficiencia de 10 empleados en dos momentos del año.
  La variable respuesta $Y$ es la evaluación de cada trabajador a final de año, y la variable
  explicativa $X$, la evaluación a mitad de año. Para cada valor de $X$ existe una
  distribución de probabilidad de $Y$, y la \emph{curva de regresión} une las medias de esas
  distribuciones.

  \begin{center}
    \begin{tikzpicture}[x=0.09cm, y=0.06cm]
      % ejes
      \draw[-{Stealth}] (40,40) -- (105,40) node[below] {$X$};
      \draw[-{Stealth}] (40,40) -- (40,100) node[left] {$Y$};
      \foreach \t in {50,70,90} {
        \draw (\t,40) -- (\t,39) node[below, font=\scriptsize] {\t};
        \draw (40,\t) -- (39,\t) node[left, font=\scriptsize] {\t};
      }
      \node[below, font=\small] at (72,32) {Evaluación a mitad de año};
      \node[rotate=90, font=\small] at (29,70) {Evaluación final};
      % curva de regresión m(x)
      \draw[NavyBlue, thick, domain=45:100, samples=40]
        plot (\x, {50+0.9*(\x-45)-0.006*(\x-45)^2});
      \node[NavyBlue, font=\scriptsize, anchor=west] at (100,80.8) {curva de regresión};
      % densidades condicionadas (tumbadas)
      \foreach \xx in {55,72,90} {
        \pgfmathsetmacro{\mm}{50+0.9*(\xx-45)-0.006*(\xx-45)^2}
        \draw[gray] (\xx,{\mm-15}) -- (\xx,{\mm+15});
        \draw[colRes, thick, domain=-15:15, samples=40]
          plot ({\xx+9*exp(-(\x)^2/40)}, {\mm+\x});
        \fill[NavyBlue] (\xx,\mm) circle (0.5pt);
      }
      \node[colRes, font=\scriptsize, align=center] at (82,100) {distribución de $Y$\\ para cada $X$};
    \end{tikzpicture}
  \end{center}
\end{ejemplo}

\paragraph{Filosofía.}
\begin{itemize}
  \item La variabilidad de la respuesta $Y$ depende de muchas causas, quizá infinitas.
  \item De ellas, unas pocas, $X_1,\ldots,X_k$, son \emph{causas importantes} que pueden
    observarse y controlarse.
  \item El resto, $X_{k+1},X_{k+2},\ldots$, son muchas otras causas no observables,
    desconocidas o incontrolables.
\end{itemize}
Por tanto, $Y$ puede modelizarse como
\[
  Y = f(X_1,\ldots,X_k) + g(X_{k+1},X_{k+2},\ldots),
\]
donde $g(X_{k+1},X_{k+2},\ldots)=\varepsilon$ es una \textbf{perturbación aleatoria}. Así:
\[
  Y = f(X_1,\ldots,X_k) + \varepsilon .
\]

\subsection{Objetivos y decisiones del análisis de regresión}

\paragraph{Objetivos.} El análisis de regresión tiene varios objetivos:
\begin{enumerate}
  \item \textbf{Conocer el modelo}, para averiguar el tipo o forma de la relación (lineal,
    polinómica\ldots), medir su fuerza y comprender el papel y la importancia de cada variable
    explicativa.
  \item \textbf{Predecir} el valor de $Y$ en un individuo a partir de sus valores de
    $X_1,\ldots,X_k$.
  \item \textbf{Optimizar procesos}: averiguar los valores de $X_1,\ldots,X_k$ que proporcionan
    el valor óptimo de $Y$.
\end{enumerate}

\begin{ejemplo}[Usos de los modelos de regresión]
  \begin{itemize}
    \item Estimar la ganancia de peso de niños debida a varios suplementos dietéticos.
    \item Predecir el resultado académico a partir de las calificaciones de bachillerato y la
      EBAU.
    \item Estimar la cantidad de ventas en relación con los gastos de publicidad.
    \item Predecir el consumo de calefacción según las temperaturas diarias y otros factores
      climáticos.
  \end{itemize}
\end{ejemplo}

\paragraph{Decisiones importantes.}
\begin{itemize}
  \item Selección de las ``mejores'' variables explicativas.
  \item Forma de la componente sistemática de la relación: lineal, cuadrática\ldots
  \item \textbf{Alcance (\emph{scope}) del modelo}. Uno de los usos principales de la
    regresión es la predicción, por lo que hay que tener en cuenta el \emph{rango y la
    dispersión de los valores de $X$ observados} para determinar el rango de validez de las
    predicciones del modelo ajustado. Conviene tenerlo presente ya en la fase de recogida de
    datos.
\end{itemize}

\subsection{El modelo de regresión lineal}

Si $Y=f(X_1,\ldots,X_k)+\varepsilon$, la regresión busca la función $f$ que mejor se ajusta a
los datos. Cuando $f$ es lineal, el modelo es
\[
  Y = \beta_0 + \beta_1X_1 + \beta_2X_2 + \cdots + \beta_kX_k + \varepsilon,
\]
y el objetivo se reduce a \textbf{estimar los coeficientes} $\beta_i$. La asignatura se centra
en los modelos de regresión lineal por varios motivos:
\begin{itemize}
  \item son simples y fáciles de manejar;
  \item la mayoría de los modelos pueden linealizarse mediante transformaciones;
  \item casi cualquier función se puede aproximar localmente por funciones lineales.
\end{itemize}

\begin{ejemplo}[Tipos de regresión lineal][ej:t2-tipos]
  \begin{itemize}
    \item \textbf{Regresión simple} (un solo predictor):
      \begin{itemize}
        \item Número de accidentes de tráfico mortales frente al número de habitantes de cada
          estado: la nube de puntos es claramente creciente y aproximadamente lineal.
        \item La longitud de un caimán es fácil de medir con fotografías aéreas, pero su peso
          es mucho más difícil de estimar, por lo que se quiere predecir el peso a partir de
          la longitud. El peso crece de forma \emph{no lineal} con la longitud, pero de forma
          aproximadamente lineal con el cubo de la longitud; el modelo
          $Y=\beta_0+\beta_1X^3+\varepsilon$ es lineal en los parámetros
          (\cref{def:modelo-lineal}). Es un ejemplo de relación \emph{linealizable}.
      \end{itemize}
    \item \textbf{Regresión lineal múltiple}: consumo de un vehículo (millas por galón) frente
      a su potencia y su peso. La superficie ajustada es un plano.
    \item \textbf{Regresión lineal a trozos}: descripción de la morfología de las glándulas de
      Meibomio mediante una poligonal.
  \end{itemize}
\end{ejemplo}
